\scanpage{567}
\mspage{478}
\deleted{de la \inserted{sous-}catégorie \deleted{\unsure{str. pleine}} ouverte de $M$
formée des $x\in\operatorname{Ob}M$ tels que $\operatorname{Hom}_M(x,z)\ne\varnothing$.}

\deleted{En effet, toute \deleted{\unsure{ouverte est}} sous-catégorie ouverte
\inserted{$\unsure{\underline T|\overline{\underline I}|\overline{\underline J}}$} est une réunion des sous-catégories ouvertes
de la forme $M_z\unsure{.}$ (pour $z\in\operatorname{Ob}\underline T$). Comme
l'ensemble des \deleted{\unsure{ces}} sous-catégories ouvertes
de $M$ est petit (celui des sous-catégories
fermées l'étant), il s'ensuit que toute
sous-catégorie ouverte est une réunion
d'une famille \uline{petite} de sous-catégories
$M_z$. Si ces dernières sont accessibles
dans $M$, il en est donc de $=$ de leurs
réunions en vertu de 4.13\ldots, d'où le
corollaire.}

\deleted{Démonstration de 4.16.7.1. Je vais d'abord
prouver b). Notons le

\uline{Lemme} 4.16.7.3. Soit \deleted{\unsure{$C$}} $f:M$\inserted{$\longrightarrow N$} \unsure{une} \deleted{catégorie,}
\deleted{\unsure{$N$ une sous}} un foncteur pl. fidèle entre
deux catégories. Alors l'application
\[
\alpha:N'\longmapsto f^{-1}(M'):\unsure{\mathcal F|F}(N)\mathrel{\unsure{-|\longrightarrow}}\mathcal F(M)
\]
est surjective, où $\mathcal F(N)$ est \unsure{l'ens.} des
\deleted{\unsure{parties}}\inserted{ss-cat.} fermées \deleted{\unsure{du}}de $\deleted{\unsure{M}}N$. De façon plus
précise, considérons l'application \deleted{\unsure{$\beta:\mathcal F(M)\longrightarrow\mathcal F(N)$}, $\unsure{M'\mapsto N'}$, où $\unsure{N'=f(M')|N'=f_!M'}$ \inserted{$\unsure{f_!M'|f(M')}$} \unsure{ensemble} $\unsure{f(\operatorname{Ob}M)}$}}
% archive sheet number: 566
